\section{Introduction}

Modern medication prescribing is a complex clinical decision-making process that requires balancing therapeutic efficacy, patient-specific constraints, dosage considerations, and potential adverse drug interactions. Physicians must select from thousands of available drug options while accounting for comorbidities, demographics, and safety risks. This cognitive burden often leads to variability in prescribing patterns and increases the risk of suboptimal or unsafe medication combinations. As electronic health records (EHRs) have become widely available, researchers have increasingly explored data-driven approaches to assist in medication recommendation and clinical decision support.

Deep learning (DL) methods have demonstrated strong capability in modeling complex, high-dimensional healthcare data. Early work such as \textit{Doctor AI} \cite{choi2016doctorai} employed recurrent neural networks (RNNs) over longitudinal EHR data to predict future diagnoses and medication categories. By learning temporal patient embeddings from sequences of visits, these models captured disease progression patterns and improved multi-label prediction performance. Similarly, \textit{RETAIN} \cite{choi2016retain} introduced a reverse-time attention mechanism that enhanced interpretability while maintaining strong predictive accuracy. By assigning attention weights to past visits and clinical variables, RETAIN allowed clinicians to better understand model reasoning, addressing a key concern in deploying deep learning systems in healthcare.

Beyond sequence modeling, recent research has incorporated graph-based representations and external medical knowledge to improve safety and robustness. \textit{GAMENet} \cite{shang2019gamenet} integrated patient history encoding with a Graph Convolutional Network (GCN) operating over drug--drug interaction (DDI) graphs. By explicitly modeling adverse drug interactions, GAMENet reduced unsafe medication combinations while preserving recommendation accuracy. Other approaches, such as \textit{SafeDrug} \cite{lee2021safedrug}, leveraged molecular graph encoders to capture chemical substructures and improve DDI-aware predictions. These advances illustrate how deep architectures can combine temporal modeling, graph neural networks, and domain knowledge to enhance both effectiveness and safety in medication recommendation systems.

However, most existing approaches assume rich, longitudinal EHR data typically available in tertiary-care or inpatient hospital settings. In contrast, outpatient and primary-care environments often involve sparse data, limited historical visits, and short free-text symptom descriptions. Many encounters consist of a single visit with minimal structured history, which reduces the effectiveness of purely sequential models. Furthermore, outpatient medication guidance must explicitly prioritize safety, ensuring that predicted drug combinations minimize harmful interactions even when historical context is limited.

In this work, we address these challenges by conducting a focused survey of deep learning-based medication recommendation systems, with particular emphasis on safety modeling, interpretability, and applicability to data-sparse outpatient settings. Based on identified gaps, we propose a hybrid deep learning architecture that integrates attention-based symptom encoding, structured feature fusion, and DDI-aware safety filtering. Our objective is to design a medication recommendation framework that balances predictive performance, safety constraints, and interpretability, making it suitable for primary-care clinical decision support.

